% §1 Introduction — Target 1.5pp

\section{Introduction}
\label{sec:intro}

Bitcoin's custody model places the entire burden of key security on the
asset holder: a compromised private key results in irreversible fund loss
with no protocol-level recourse. Covenant-based vaults address this by
interposing a time-delayed spending path between key compromise and fund
movement, so that an unauthorized hot-key spend can be intercepted by a
watchtower during a configurable delay window and redirected to cold
storage~\cite{SHMB20,MES16}. We focus on four Bitcoin Script
extension proposals with published BIPs, reference implementations,
and vault-specific threat discussion in the 2023--2026 activation
debate: \opctv{}~\cite{BIP119}, \opccv{}~\cite{BIP443}, \opvault{}~\cite{BIP345}
(withdrawn~\cite{OB25}, included for comparison), and the
\opcat{}+\opcsfs{} composition~\cite{BIP347,BIP348}. Other
covenant-capable proposals (ANYPREVOUT, TXHASH, TLUV) are beyond our
scope. Practitioners must choose among the four, yet no peer-reviewed
cross-design comparison exists.

The four proposals differ in how they introspect the spending
transaction, how they authorize recovery, and how they pay fees. These
differences are not cosmetic: each choice opens a distinct
vulnerability surface---anchor-based fee models admit descendant-chain
pinning (class \classfee{}), partial-withdrawal recovery admits a
per-UTXO recovery-cost tail whose ordering between \opccv{} and
\opvault{} inverts with watchtower batching (class \classscale{}),
keyless recovery admits permissionless griefing (class \classgrief{}),
destination-substitutable withdrawal admits hot-key theft
(class \classhot{}), and unbound recovery leaves admit
cold-key theft (class \classcold{}). No single design minimises
defender loss across all classes and all fee rates. The
activation debate~\cite{CTVCSFS25,Poinsot25,Towns25,OB25} has proceeded
without quantitative input on these tradeoffs; the literature contains
only qualitative matrices~\cite{covenants_info} and a single
back-of-envelope fee-pinning estimate~\cite{Har24} based on assumed
rather than measured transaction sizes.

This paper fills that gap. We build all four Bitcoin vault
implementations behind a uniform lifecycle interface and run 15
experiments on regtest measuring transaction sizes, attack costs,
and per-design functional differences. We additionally develop Alloy
bounded model checking specifications for structural properties (fund
conservation, recovery liveness, no-extraction), and derive
closed-form expressions for dust-threshold fee rates, pinning costs,
and batching amortisation. We implement a fifth design---a Simplicity
vault on Elements---as a cross-substrate reference point; because
Elements uses federated rather than proof-of-work consensus, the
Simplicity results are presented in Appendix~\ref{app:simplicity} and
are excluded from the formal statements below. Simplicity serves as
a constructive existence proof that the dominant corner of the
Bitcoin design space is occupiable on some substrate, even
while no Bitcoin proposal occupies it.

\paragraph{Contributions.}
\begin{enumerate}
  \item A \textbf{structural design-space decomposition}: covenant
    vault designs are characterised by seven canonical axes
    (\axauth{}, \axrevault{}, \axfee{}, \axwdbind{}, \axrecbind{},
    \axopsurf{}, \axintro{}; Section~\ref{sec:bg:axes}), with six
    derived attack classes \classfee{}--\classmode{}. Each class
    is a deterministic consequence of a specific axis-value
    (Sections~\ref{sec:imposs:class_fee}--\ref{sec:imposs:class_mode}),
    susceptibility and immunity both mechanical readings of the
    position table. No proposed Bitcoin extension occupies
    the dominant corner (Section~\ref{sec:imposs:design_space}).
  \item A \textbf{Griefing--Safety Incompatibility} proposition: no
    covenant vault simultaneously achieves permissionless recovery and
    griefing resistance (Section~\ref{sec:imposs:class_grief}).
  \item A \textbf{per-threat defender-loss characterisation}
    $\mathcal{L}(\mathrm{TM}, D, V, r, b)$ parameterised by threat
    model, design, vault value, fee rate, and defender batching
    strategy (Section~\ref{sec:imposs:sensitivity}).
  \item \textbf{Batched-defender measurements and a recovery-cost
    ordering-flip finding.} When the watchtower recovers many
    triggered vault UTXOs in one transaction instead of one at a
    time, its per-input cost drops sharply, raising the fee rate
    at which the per-UTXO recovery-cost tail of class \classscale{}
    saturates within a single block by $1.8\times$ for \opccv{} and
    $2.7\times$ for \opvault{}. The safer of the two designs under
    \classscale{} therefore depends on whether the watchtower
    batches: \opccv{} is safer without batching, \opvault{} is safer
    with it (Section~\ref{sec:empirical:batching}).
  \item \textbf{Axis-keyed deployment guidance} keyed to fee
    regime, batching strategy, and adversary profile, derived as
    counter-factual readings of the class rationality blocks
    (Section~\ref{sec:deployment}).
\end{enumerate}

Contributions are derived from the design-space decomposition of
Section~\ref{sec:bg:axes} and referenced uniformly throughout
Sections~\ref{sec:impossibility}--\ref{sec:deployment}. Every
attack-feasibility claim in this work is evaluated under the
Kerckhoffs threat model of Section~\ref{sec:bg:lifecycle}.

\paragraph{Scope.} Our formal results range over the four proposed
Bitcoin Script extensions only. We implement and discuss
Simplicity~\cite{OC17}, an alternative script primitive deployed on
the Elements federated sidechain~\cite{DPNS16}, as a cross-substrate
reference point (Appendix~\ref{app:simplicity}). Elements'
functionary-federated consensus differs fundamentally from Bitcoin's
proof-of-work, and the threats we analyse do not transfer identically
across that substrate boundary. Protocol-layer constructs such as Ark
and BitVM are also excluded: both operate above the opcode level and
require modelling interactive rounds outside our single-transaction
covenant framework.

\paragraph{Roadmap.}
Section~\ref{sec:background} summarises the four Script extensions,
the vault lifecycle, and the 11-element threat taxonomy.
Section~\ref{sec:impossibility} presents the structural results.
Section~\ref{sec:empirical} reports the regtest measurements.
Section~\ref{sec:deployment} derives deployment guidance from the
per-threat analysis. Section~\ref{sec:related} positions this work
against prior vault research and Section~\ref{sec:discussion}
discusses limitations and open problems.

\providecommand{\todo}[2][]{\textbf{[TODO: #2]}}
